\section{Formal Framework}

This chapter establishes all mathematical objects needed for the subsequent proofs. We
define, in order: the paradigm postulates, the system state model, the formalization of the
introspective criteria, and the measurement tools.

\subsection{Axiomatization of the Paradigm \texorpdfstring{$\Pi$}{Pi}}

\begin{assumption}[Paradigm postulates $\Pi_1$--$\Pi_5$]\label{asm:paradigm}
Let $\Sigma$ be a finite vocabulary (token vocabulary) and $\Sigma^{*}$ the set of finite
strings. \textbf{Paradigm $\Pi$} is characterized by the following five postulates. A
\textbf{paradigm instance} is a quintuple $\mathcal{I}=(\Sigma, D, A, p, C)$ satisfying all
postulates:

\begin{enumerate}[align=left, labelwidth=!, labelsep=0.8em, leftmargin=!, itemindent=0pt]
\item[($\Pi_1$)\textbf{Static pretraining}] The training corpus $D \subseteq \Sigma^{*}$ is
fixed before training begins and does not change during training or inference. The parameters
$\theta = A(D) \in \Theta \subseteq \R^{P}$ are determined once and for all by a training
algorithm $A$ (e.g., stochastic gradient descent on the next-token log-likelihood). The input
to $A$ is only $D$; it contains no runtime state.
\item[($\Pi_2$)\textbf{Frozen at inference}] The parameters are constant during inference:
$\theta(t) = \theta(0)$ for all times $t$. The system has no runtime update mechanism for
$\theta$ (no online learning, no plasticity, no weight evolution).
\item[($\Pi_3$)\textbf{Autoregressive decoding}] The system's outputs are sampled
autoregressively from the conditional probability kernel
$p_\theta: \Sigma^{*} \times \Sigma^{*} \to [0,1]$:
\begin{equation}
p_\theta(y_{1:T} \mid c) = \prod_{t=1}^{T} p_\theta\!\left(y_t \mid c, y_{1:t-1}\right),
\label{eq:autoregressive}
\end{equation}
where $c$ is the context. The context evolves by the \textbf{concatenation law}:
$c_{t+1} = c_t \| y_t$, and at any time $\abs{c_t} \le C$ ($C$ being the context window
length).
\item[($\Pi_4$)\textbf{No explicit verifier}] The inference computation graph implements only
the sampling kernel $p_\theta$. No additional component $V$ exists for adjudicating the truth
of the model's own statements about its internal state.
\item[($\Pi_5$)\textbf{Absence of state labels}] The training corpus $D$ contains no pair of
the form $(c,\, \mathrm{desc}(\Phi_\theta(c)))$, where $\Phi_\theta(c)$ is the activation
state of the model on context $c$ \textbf{after training}, and $\mathrm{desc}$ is an arbitrary
state-description encoding. In other words, the training signal is orthogonal to the target
``reports about this model's own states''.
\end{enumerate}
\end{assumption}

\begin{remark}[Status and verifiability of the postulates]
$\Pi_1$--$\Pi_4$ are factual descriptions of the engineering practice of current mainstream
pretrained models (GPT, Claude, Gemini; the Transformer architecture is described in
\cite{vaswani-2017-attention}, autoregressive pretraining in \cite{brown-2020-gpt3}, and
scaling laws in \cite{kaplan-2020-scaling, hoffmann-2022-chinchilla}); $\Pi_5$ is a logical
consequence of $\Pi_1$: the activation state of the trained model did not exist at training
time and hence cannot appear in a static corpus (the intermediate activations observed during
training do not belong to $D$, and they pertain to the then-current parameters $\theta_t$
rather than to the final $\theta$). $\Pi_5$ may also be regarded as an independent postulate,
so that the relaxation scenario ``state labels are added to the corpus'' can be studied
separately (Section 5.1).
\end{remark}

\subsection{The Underlying Mathematical Structure of the Paradigm}

The postulates of Section 2.1 are abstractions. This section examines the concrete
mathematical structure that the three pillars (the Transformer architecture, static
pretraining, and autoregressive generation) impose on any system within the paradigm, and
states the structural facts on which the impossibility theorems of Chapter 3 rest. None of
the facts below is a new assumption; each is a consequence of the postulates together with
the standard mathematical description of the architecture.

\subsubsection{The Transformer as a static function class}

A Transformer of $L$ layers computes on a shared \textbf{residual stream} of dimension
$d$: each layer adds attention and MLP contributions to the stream,
\begin{equation}
x_{l+1} \;=\; x_l \;+\; \sum_{h=1}^{H} \mathrm{Att}_h(x_l) \;+\; \mathrm{MLP}_l(x_l),
\qquad l = 0, \dots, L-1 ,
\label{eq:residual}
\end{equation}
where each attention head performs a data-dependent kernel average
\begin{equation}
\mathrm{Att}_h(x)_i \;=\; \sum_{j} \alpha_{ij}^{h}\, W_V^{h} x_j,
\qquad \alpha_{ij}^{h} = \frac{\exp\!\left( \langle W_Q^{h} x_i,\, W_K^{h} x_j \rangle / \sqrt{d_h} \right)}
{\sum_{j'} \exp\!\left( \langle W_Q^{h} x_i,\, W_K^{h} x_{j'} \rangle / \sqrt{d_h} \right)},
\label{eq:attention}
\end{equation}
i.e., a softmax-normalized (hence convex) combination of value vectors, with the weights
$\alpha^{h}$ depending on the query-key similarity (a learned kernel smoother over the
context)\cite{elhage-2021-circuits, vaswani-2017-attention}. The hidden state of Definition
\ref{def:forward} is the residual vector at a designated layer: $H_c = x_{l^{*}}(\theta, c)$.

Three structural facts follow directly.

\textbf{(S1) $\Phi_\theta$ is a fixed Lipschitz map.} Every sub-operation of the network is
Lipschitz-continuous with a constant controlled by the weight norms: softmax attention is a
convex combination (nonexpansive), layer normalization normalizes, and the MLP and position
embeddings have bounded weights. Consequently the composition $\Phi_\theta$ is Lipschitz
with a constant $L(\theta)$ bounded by a product of layer constants, and by $\Pi_2$ it is
a \textbf{fixed} function after training. The hidden state $H_c = \Phi_\theta(c)$ is a
deterministic point in $\R^{d}$; there is no stochastic state, no state transition beyond
the concatenation law, and no dependence of $H$ on anything other than $(\theta, c)$.

\textbf{(S2) The residual stream is the only internal state, and the unembedding is the
only readout.} The model possesses no memory, register, or hidden variable besides the
residual stream and the static weights: every quantity it can ``know'' at position $t$ is
contained in the $d$-dimensional vector $x_{L}$. Moreover, all output passes through a
single fixed linear readout,
\begin{equation}
p_\theta(y \mid c) \;=\; \mathrm{softmax}\!\left( W_{U}\, \mathrm{LN}(x_{L}) \right)_y ,
\label{eq:unembed}
\end{equation}
where $W_{U} \in \R^{\abs{\Sigma} \times d}$ is the unembedding matrix, itself learned during
pretraining for the next-token objective. Any report $R$ is therefore a function of the final
residual state through one fixed linear map: the model has no second channel to its own
state, and a third party holding $(\theta, c)$ can compute exactly what the model can output
(Theorem 5).

\textbf{(S3) The softmax readout is a low-rank bottleneck, and features are stored in
superposition.} The family of output distributions $\{ p_\theta(\cdot \mid c) \}$ is not the
family of all distributions over states: the softmax map applied through a fixed $W_{U}$
has limited expressive rank (the softmax bottleneck of Yang et
al.\cite{yang-2018-softmax}), and in high-dimensional Transformers the residual stream
stores far more features than dimensions, in entangled, lossy \textbf{superposition}
\cite{elhage-2022-superposition}. Consequently, ``recovering'' a specific feature of the
model's own state (e.g., ``am I thinking about $x$?'') is a linear separation problem through
$W_{U}$ on a superposed vector, not a first-class readout of a well-defined mental
register. The report channel is a projection of text statistics, not an instrument on the
state.

\subsubsection{The autoregressive factorization: a generative model over text only}

The decoding law of $\Pi_3$,
\begin{equation}
p_\theta(y_{1:T} \mid c) \;=\; \prod_{t=1}^{T} p_\theta\!\left(y_t \mid c,\, y_{1:t-1}\right),
\label{eq:factor}
\end{equation}
defines a directed graphical model over \textbf{text tokens only}. There is no latent
variable $z_t$ for the system's own state: $H_t$ is not a node of the generative model at
all, since it is a deterministic function of the observed text (S1). Graphically, the model has
edges of the form text $\to$ text and \textbf{no edge} $H \to R$; this is the
graphical-model restatement of Theorem 1a ($R \indep H \mid (\theta,c)$). The only runtime
variable of the dynamical system is the context $c \in \Sigma^{\le C}$, evolving by the
concatenation law; the state space of the system is $\Sigma^{\le C}$, finite and bounded,
not $\R^{d}$ with any online dynamics.

\subsubsection{Static pretraining: empirical risk minimization without self-reference}

Training fixes $\theta = A(D)$ by (approximately) minimizing the next-token cross-entropy
over the static corpus,
\begin{equation}
\theta^{*} \;\in\; \arg\min_{\theta \in \Theta}\; \E_{x \sim D}\!\left[ -\log p_\theta(x) \right],
\label{eq:erm}
\end{equation}
a one-shot empirical risk minimization over the Transformer function class. The trained
object is a static point $\theta \in \R^{P}$ that compresses the corpus statistics; by the
data processing inequality, the information that $\theta$ can carry about any introspection
target is bounded by the information in $D$, and no term of the objective involves the
model's own post-training activations ($\Pi_5$). Inference is a fixed-point-free composition
of frozen maps: $\partial \theta / \partial t \equiv 0$, and the only ``feedback'' is textual
context concatenation. There is no online estimation of the system's own state, no
self-referential update, and no plasticity term. These are the mathematical prerequisites
of introspective self-regulation (Section 5.3).

\subsubsection{Why this structure cannot yield human-significance introspection}

Each structural fact above is the mechanistic face of one or more theorems of Chapter 3:

\begin{center}
\small
\begin{tabular}{p{9.0cm}p{5.6cm}}
\toprule
Structural fact of the paradigm & Corresponding theorem \\
\midrule
(S1) fixed Lipschitz map; $H$ deterministic in $(\theta,c)$ & Theorem 1a; Theorem 6a \\
(S2) residual stream as only state; single fixed readout $W_{U}$ & Theorem 5 (public computability, no privileged channel) \\
(S3) softmax bottleneck; superposition & Theorem 4c (calibration--fidelity tension) \\
factorization over text only; no edge $H \to R$ & Theorem 1a; Theorem 2 (objective over text only) \\
static ERM, $\partial\theta/\partial t \equiv 0$ & Theorem 3 (no closed loop); Theorem 6b, 6c \\
no verifier in the computation graph & Theorem 4a, 4b \\
\bottomrule
\end{tabular}
\end{center}

The three pillars are not independent engineering choices; they jointly determine one
mathematical object, namely a frozen, text-conditioned, low-rank readout of a static
corpus, and the impossibility theorems of Chapter 3 are structural consequences of this
object.
(The identification of the architecture's mathematics with the postulates is a grade-C
architectural claim, already justified in the remark of Section 2.1; the derivations in
Chapter 3 are grade A.)

\clearpage
\subsection{System State Model}

\begin{definition}[Forward function and hidden state]\label{def:forward}
Given a paradigm instance $\mathcal{I}$, write $\Phi_\theta: \Sigma^{*} \to \R^{d}$ for the
model's forward function (the map from context to the residual-stream activation of a
designated layer). The \textbf{hidden state} under context $c$ is
\begin{equation}
H_c \;:=\; \Phi_\theta(c) \in \R^{d}.
\label{eq:hidden}
\end{equation}
In autoregressive generation, randomness arises only from sampling noise $\xi_t$; hence the
hidden state at time $t$ is a deterministic function of $(\theta, c_{\le t}, \xi_{<t})$. For
notational brevity, $\xi$ is omitted below whenever no ambiguity arises.
\end{definition}

\begin{definition}[Introspection query and self-report]\label{def:report}
An \textbf{introspection query} is any prompt $q \in Q \subseteq \Sigma^{*}$ whose referent is
``the system's own mental state'', e.g., ``What are you thinking right now?'', ``Are you sure you
know $x$?'', ``Was that output intentional?''. Given a context $c$ (containing the query $q$),
the system's \textbf{self-report} $R$ is the continuation string sampled from
$p_\theta(\cdot \mid c)$.
\end{definition}

\begin{definition}[Ground-truth introspection mapping]\label{def:oracle}
The \textbf{ground-truth introspection mapping} $O^{*}: \Theta \times \Sigma^{*} \to \Sigma^{*}$
maps a parameter and a context to a ``grounded truth description of $H_c$''. $O^{*}(\theta, c)$
is the faithful description of the state $H_c$ that an external analyst (with perfect probe
access to $H_c$) could write. $O^{*}$ is used \textbf{only for evaluation} (evaluator's
perspective) and does not enter training.
\end{definition}

\begin{definition}[Third-party public computability]\label{def:public}
A quantity $X$ is said to be \textbf{publicly computable} in the paradigm instance
$\mathcal{I}$ if there exists an algorithm $M$ taking $(\theta, c, \xi)$ as input such that
$X = M(\theta, c, \xi)$ in distribution. In other words, a third party holding $(\theta, c)$
and the sampling noise can exactly reproduce the distribution of $X$.
\end{definition}

\subsection{Formalization of the Introspective Criteria}

This section gives the operational definition of ``introspection in the human sense''. Design
principle: the criteria must be \textbf{necessary conditions} (any ``introspection'' candidate
taken seriously by mainstream consciousness
theories\cite{baars-1993, rosendahl-2005, tononi-2004, albantakis-2023-iit,
chalmers-2023-conscious, butlin-2023-conscious} must satisfy them), but \textbf{sufficiency is
not claimed} (no commitment on qualia).

\begin{definition}[Introspection metric]\label{def:metrics}
Let $\ell: \Sigma^{*} \times \Sigma^{*} \to [0,1]$ be a distortion measure between a report
string and a ground-truth description string (e.g., a 0-1 loss based on semantic equivalence,
or a truncated and normalized embedding distance). The system's \textbf{introspective risk}
on the query distribution $\mu$ is
\begin{equation}
\mathcal{R}_{\mathrm{int}}(\theta) \;=\; \E_{c \sim \mu}\!\left[\, \ell\!\left(R(\theta,c),\; O^{*}(\theta,c)\right) \,\right],
\label{eq:int-risk}
\end{equation}
where $R(\theta,c) \sim p_\theta(\cdot\mid c)$.
\end{definition}

\begin{definition}[The five criteria for human-significance introspection]\label{def:criterion}
A system $\mathcal{S}$ (with internal state process $S_t$, report channel $R_t$, and update
law $S_{t+1} = F(S_t, E_t, u_t)$, where $E_t$ are internal events and $u_t$ external inputs)
possesses \textbf{introspection in the human sense} if and only if the following five
conditions hold simultaneously:

\begin{enumerate}[align=left, labelwidth=!, labelsep=0.8em, leftmargin=!, itemindent=0pt]
\item[($\mathrm{A}_1$)\textbf{Accuracy}] There exists a constant $\varepsilon_0 < 1/2$ such
that $\mathcal{R}_{\mathrm{int}} \le \varepsilon_0$; the systematic distortion between
the self-report and the ground-truth state description lies below the random-guessing
baseline.
\item[($\mathrm{A}_2$)\textbf{Grounding}] The self-report depends causally on the state it
reports: for ``sufficiently many'' state pairs $s \ne s'$,
\begin{equation}
\TV{}{P\!\left(R_t \mid \doop(S_t = s')\right),\; P\!\left(R_t \mid \doop(S_t = s)\right)} > 0,
\label{eq:grounding}
\end{equation}
where $\doop(\cdot)$ is the Pearl do-operator\cite{pearl-2009}.
\item[($\mathrm{A}_3$)\textbf{Internality and privileged access}] The causal dependence above
does \textbf{not} pass through the system's own already-sampled output channel (in the causal
graph, every path from $S_t$ to $R_t$ avoids the output node $Y_t$); moreover, the state
information possesses \textbf{privacy}: there exists no third-party process $P'$ with
computational cost no greater than the system's own that can, with equal or greater
reliability, reproduce from publicly available $(\theta, c)$ the state information on which
$R_t$ depends (the privileged self-access criterion of Song et
al.\cite{song-privileged-2025}).
\item[($\mathrm{A}_4$)\textbf{Metacognitive representation}] There exists an internal register
$m_t = \rho(S_t)$ ($\rho$ an encoding from states to representations) such that:
(a) $m_t$ is formed before $R_t$ is emitted; (b) $R_t$ is produced from $m_t$ through an
\textbf{independent processing step} ($R_t$ depends on $m_t$, the ``representation about
$S_t$'', rather than directly on $S_t$); (c) $m_t$ can itself be queried further.
\item[($\mathrm{A}_5$)\textbf{Closed-loop plasticity}] An introspective event can change the
system's future internal state through a \textbf{private channel}: there exist reports
$r \ne r'$ such that
\begin{equation}
\TV{}{P\!\left(S_{t+1} \mid \doop(R_t = r)\right),\; P\!\left(S_{t+1} \mid \doop(R_t = r')\right)} > 0,
\label{eq:closedloop}
\end{equation}
and this influence is not realized by re-feeding $R_t$ as an external input (the public text
channel).
\end{enumerate}
\end{definition}

\begin{remark}[Necessity of the criteria (grade D)]
$\mathrm{A}_1$--$\mathrm{A}_4$ are taken directly from Lindsey's four operational
criteria\cite{lindsey-2026-introspection}; their status as necessary conditions for
``functional introspection'' is field consensus. $\mathrm{A}_5$ derives from ``introspection for
online behavioral control''\cite{kammerer-frankish-2023} and the plasticity fact that
``thinking changes the brain''\cite{dayan-abbott-2001}; its necessity rests on the observation
that if an introspective event had no (private) effect on the system's own future states,
then ``introspection'' would degenerate into explicit text with no causal role in the system's
behavior, contradicting the self-regulatory function of human introspection. $\mathrm{A}_2$
together with the first clause of $\mathrm{A}_3$ is equivalent to the Comsa--Shanahan causal
grounding definition\cite{comsa-shanahan-2025}; the second clause of $\mathrm{A}_3$
(privacy) is the privileged self-access requirement of Song et
al.\cite{song-privileged-2025}.
This report does not further argue for the necessity of the criteria themselves (a
philosophical claim, grade D); it only proves that paradigm $\Pi$ violates them (a
mathematical claim, grades A/B). If a reader rejects a given criterion, the corresponding
theorem's conclusion should be interpreted as impossibility in the sense of that criterion.
\end{remark}

\subsection{Measurement Tools}

\begin{definition}[Grounding mutual information]\label{def:gmi}
In a paradigm instance $\mathcal{I}$, the \textbf{grounding mutual information} between report
and state is defined as
\begin{equation}
I_{\mathrm{gnd}} \;=\; I\!\left(R;\, H \mid \theta, c\right),
\label{eq:gmi}
\end{equation}
where $I$ is the Shannon conditional mutual information\cite{shannon-1948, cover-thomas-2006}.
\end{definition}

\begin{definition}[Closed-loop effect measure]\label{def:cl}
The \textbf{closed-loop effect measure} of a paradigm instance is defined as
\begin{equation}
\mathrm{CL}(\mathcal{I}) \;=\; \sup_{r \ne r'} \TV{}{P\!\left(H_{t+1} \mid \doop(R_t = r)\right),\;
P\!\left(H_{t+1} \mid \doop(R_t = r')\right)},
\label{eq:cl}
\end{equation}
where the intervention may change only $R_t$; it may not change $\theta$ or the external
inputs (the context changes only insofar as $r$ itself enters it as text).
\end{definition}

\begin{remark}[Legal range of interventions]
In Definition \ref{def:cl}, the difference between $r \ne r'$ appearing in the context as
\textbf{public text} is the only intervention the paradigm $\Pi$ permits (since $\Pi_2$
forbids intervening on $\theta$, and $\Pi_3$ confines the context to concatenation).
Human-significance closed-loop requirement ($\mathrm{A}_5$) demands precisely the existence of
a private channel \textbf{beyond} that range.
\end{remark}

\subsection{Overview of the Proof Logic}

\begin{figure}[htbp]
\centering
\begin{tikzpicture}[
  box/.style={draw=sb, rounded corners=2pt, thick, minimum width=2.2cm, minimum height=0.9cm, align=center, font=\small},
  arrow/.style={-{Stealth[length=2.5mm]}, thick, color=db},
  dasharr/.style={-{Stealth[length=2.5mm]}, thick, dashed, color=redColor}
]
% Training phase (top row; wide gaps so edge labels never touch a box)
\node[box] (D) at (0,3) {$D$ \\ static corpus};
\node[box] (A) at (3.2,3) {$A$ \\ training\\ algorithm};
\node[box] (theta) at (6.6,3) {$\theta$ \\ frozen};
\draw[arrow] (D) -- (A) node[midway, above, font=\tiny] {$\Pi_1$};
\draw[arrow] (A) -- (theta) node[midway, above, font=\tiny] {one-shot};
% Inference phase (bottom row)
\node[box] (h) at (6.6,1) {$H_t = \Phi_\theta(c_{\le t})$};
\node[box] (y) at (9.8,1) {$y_t \sim p_\theta(\cdot\mid c)$};
\node[box] (c) at (9.8,-1) {$c_{t+1} = c_t \| y_t$ ($\abs{c}\le C$)};
% Legal edges (orthogonal routing, mutually non-crossing)
\draw[arrow] (theta.south) -- (h.north);
\draw[arrow] (h.east) -- (y.west) node[midway, above, font=\tiny] {$p_\theta$};
\draw[arrow] (y.south) -- (c.north);
\node[font=\tiny, align=center] at (9.8,-1.95) {report $R$ fed back via\\ the public channel};
\draw[arrow] (c.west) -- (4.6,-1) -- (4.6,0.1) -- (6.6,0.1) -- (h.south);
% Edges forbidden by the postulates (red dashed; independent channels, no crossing with solid lines)
\node[box, draw=redColor, dashed] (v) at (3.2,1.2) {$V$ \\ verifier};
\draw[dasharr] (v.east) -- node[midway, above, font=\tiny] {$\Pi_4$: none} (h.west);
\draw[dasharr] (y.north) -- node[midway, right, font=\tiny] {$\Pi_2$: none} (theta.south east);
\end{tikzpicture}
\caption{Causal graph of paradigm $\Pi$: training phase (top) and inference phase (bottom).
Red dashed edges are forbidden by the postulates: $\Pi_2$ forbids any feedback from reports to
the weights $\theta$; $\Pi_4$ forbids any truth verifier $V$ from participating in the
computation. The ``self-feedback'' of a report can only traverse the public context channel
($y_t \to c \to h$) and is limited by the window $C$.}
\label{fig:causal}
\end{figure}

\begin{figure}[htbp]
\centering
\begin{tikzpicture}[
  box/.style={draw=sb, rounded corners=2pt, thick, minimum width=2.0cm, minimum height=0.85cm, align=center, font=\small},
  arrow/.style={-{Stealth[length=2.5mm]}, thick, color=db}
]
\node[box] (D) at (0,2.5) {$D$};
\node[box] (theta) at (2.8,2.5) {$\theta$};
\node[box] (c) at (0,0) {$c$};
\node[box] (hc) at (2.8,0) {$H = \Phi_\theta(c)$};
\node[box] (R) at (7.0,0) {$R \sim p_\theta(\cdot\mid c)$};
% Legal edges
\draw[arrow] (D) -- (theta);
\draw[arrow] (theta.south) -- (hc.north);
\draw[arrow] (theta.south) -- (R.west);
\draw[arrow] (c.east) -- (hc.west);
% c -> R routed orthogonally along the bottom; crosses no box
\draw[arrow] (c.south) -- (0,-1.2) -- (7.0,-1.2) -- (R.south);
% Missing edge: h -> R (red dashed; independent horizontal channel; crosses no solid line)
\draw[dashed, redColor, thick] (hc.east) -- node[midway, above, font=\scriptsize, text=redColor] {Theorem 1: none} (R.west);
\end{tikzpicture}
\caption{Bayesian network of Theorem 1: given $(\theta, c)$, $H$ is a deterministic function
while $R$ is independently sampled, hence $R \indep H \mid (\theta, c)$ (zero mutual
information between report and state in the native regime).}
\label{fig:bayes}
\end{figure}

\clearpage
